\section{Introduction}
\label{sec:introduction}

Delayed feedback is intrinsic to many biological oscillators.  Signal
transmission, synaptic processing, transcription, and transport take time,
while the local units themselves often evolve on sharply separated time
scales.  In an excitable neuronal model these two features meet near the
transition from a small subthreshold response to a large voltage pulse.  The
FitzHugh--Nagumo equations provide a standard low-dimensional description of
this transition, and their coupled and spatially extended variants continue
to serve as prototypes for neuronal and other biological media; see
\citet{cebrianlacasa2024six} for a recent review.  Near a fold of the critical
manifold, the relevant orbit may follow both an attracting and a repelling
slow branch.  The parameter at which the two continuations connect is then
exceptionally sensitive: a displacement that is algebraically small in the
time-scale ratio can decide whether the response remains small or develops
into a pulse.

For planar ordinary differential equations, the geometry behind this
sensitivity is well understood.  The blow-up analysis of fold and canard
points by \citet{krupa2001extending}, together with the analysis of relaxation
oscillations and canard explosion in \citet{krupa2001relaxation}, identifies
the maximal canard as an intersection of attracting and repelling slow
manifolds and resolves the associated threshold on the appropriate inner
scale.  A delayed equation changes the mathematical object in two ways.
First, its state at a given time is a history segment rather than a point.
Second, delay pathways that are invisible after projection onto a critical
mode can still excite stable transverse modes.  Those modes may return to the
critical equation through the nonlinear geometry of the fold.  Thus a scalar
reduction can reproduce the delay operator seen directly in the critical
direction and nevertheless miss the first change in the connection
parameter.

This paper asks whether that loss of information actually occurs in a
concrete slow--fast retarded functional differential equation (RFDE).  More
precisely, suppose that two delayed networks have the same total delayed gain
and the same complete delay measure after projection onto the critical fold
mode.  Must they have the same local canard connection parameter to the first
nontrivial order?  It is important here that the projected measures, and not
only their first few moments, are identical.  We answer the question in the
negative for a two-module FitzHugh--Nagumo network with weak feedback and two
physical delays of order \(\eps^{-1/2}\).

This scaling is singular in a specific way.  The delayed feedback has gain
\(\eps K\), but its physical delay diverges as
\(\eps^{-1/2}\).  After the fold time rescaling, the two delays become the
fixed numbers \(\theta_0\) and \(\theta_1\), so the history differences
survive in the inner problem.  The setting is therefore distinct from a
small-delay expansion, in which translated histories are replaced by time
derivatives and an inertial or slow manifold can be expanded in the delay;
compare \citet{chicone2003inertial}.  It is also stronger than matching a
finite list of delay moments.  Equality of the complete critical projected
measure means that the projected delay functional agrees for every scalar
critical history.  Any threshold displacement in the family below must
therefore be produced outside that projected scalar equation.

The issue lies close to, but is not settled by, the existing theory of
delayed canards.  \citet{krupa2016explosion} studied canard explosions in DDEs
with a one-dimensional slow manifold and showed, for delayed van der Pol
dynamics, how the classical family ends at a Bogdanov--Takens bifurcation as
the delay increases.  In the neuronal FitzHugh--Nagumo setting,
\citet{krupa2016complex} described the small-delay canard regime and the
delay-induced mixed-mode oscillations, bursting, and chaotic behavior that
appear beyond it.  These works establish that delay can alter both the local
fold dynamics and the surrounding oscillatory regimes.  They do not address
two network delay organizations whose entire critical projection agrees.

A complementary infinite-dimensional viewpoint was developed by
\citet{avitabile2020local}.  Their center-manifold framework treats systems
in which an infinite-dimensional fast subsystem is coupled to finitely many
slow variables, and it gives a rigorous local reduction near bifurcations of
the layer problem, with DDEs among its applications.  This theory explains
why finite-dimensional local dynamics can persist in a broad class of
spatio-temporal problems.  The object studied below requires additional
information: a complete-history invariant graph must be controlled across a
fold region whose inner length grows as the perturbation parameter tends to
zero, and the two one-sided continuations must be matched without solving a
backward initial-value problem for the RFDE semiflow.

Classical RFDE center-manifold theory provides finite-dimensional invariant
graphs and reduced flows for each fixed regular equation
\citep{diekmannvangils1991center}; smooth parameter-dependent versions are
also available for regular fixed-delay families
\citep{bosschaert2020switching}.  A parameterization approach to delay
perturbations, based on solving invariance equations with a contraction and
then recovering higher regularity, was developed by
\citet{yang2021parameterization}.  These results motivate the history-graph
viewpoint used here, but they do not directly give estimates uniform through
the singular block \(A/\del\) or on the logarithmically growing fold tube.
We therefore prove the model-specific graph and mixed-jet estimates needed
for the connection law.

The closest scalar calculation is the recent work of
\citet{zhang2026high}.  For a van der Pol oscillator with weak delayed
feedback, they reduce the DDE on a nonlocal center manifold, expand the
resulting planar flow, and use a nonlinear time transformation to compute
high-order approximations of both the critical parameter and the critical
canard manifold.  Their result provides an important single-oscillator
comparison for the weak-feedback regime.  Our question is different even at
leading order: the directly projected delayed equation is held fixed while a
transverse delay channel is varied.  The resulting coefficient therefore
cannot be obtained from the scalar projected measure.  Its proof also needs
an invariant complete-history graph and a one-sided history connection,
rather than a scalar reduced threshold calculation; in particular, equality
must be established at the level of retained histories.

There is a separate RFDE issue behind the word ``connection.''  A current
state in \(\R^4\) does not determine whether two local pieces represent the
same RFDE orbit: their histories on the whole retained delay interval must
agree.  Conversely, the repelling piece cannot be selected by integrating an
arbitrary terminal point backward, because an RFDE does not in general
generate an invertible flow on its phase space.  Our construction selects
both one-sided pieces on a locally invariant history graph and matches them
there.  This distinction is immaterial in a planar ODE but essential here;
it is also why an observable crossing or a small voltage mismatch is not
used as the definition of the canard parameter.

We now describe the mechanism and the main result.  Let
\(\del=\sqrt\eps\), write the unfolding as \(\mu=\del^2\nu\), let
\(0<\theta_0<\theta_1\) be the delays in inner time, and write the
corresponding physical delays as \(\theta_k/\del\).  The model contains two
positive delay layers \(C_0^\eta\) and \(C_1^\eta\), with \(\eta\) in a fixed
compact subset of their positivity range and redistributing coupling between
the layers.  If \(r\) and \(\ell\) are the right and left critical vectors
normalized by \(\ell^{\mathsf T}r=1\), the construction has the exact
identities
\[
 C_0^\eta+C_1^\eta=B,
 \qquad
 \ell^{\mathsf T}\mathbb B_\eta(\dd\theta)r
 =\frac13\delta_{\theta_0}(\dd\theta)
  +\frac23\delta_{\theta_1}(\dd\theta).
\]
Both the total gain \(B\) and every atom of the critical projected delay
measure are independent of \(\eta\).  In contrast, for a transverse vector
\(q\) and the complementary projection \(P_\perp\),
\[
 P_\perp\mathbb B_\eta(\dd\theta)r
 =\eta q\bigl(\delta_{\theta_0}-\delta_{\theta_1}\bigr)(\dd\theta).
\]
The redistribution therefore forces a stable transverse voltage component
by the difference of two delayed critical histories.  That component feeds
back through the quadratic and cubic FitzHugh--Nagumo terms.  The effect is
invisible to any reduction that retains only the complete critical
projection, although it changes the local history connection.

To state the result, one local selection must be fixed.  A local
connection through the fold is chosen by an admissible preparation datum
\(\cP\).  It specifies the fixed cutoffs used to extend the reduced field,
the phase condition at the matching section, and the two tail level
conditions.  These choices define attracting and repelling one-sided traces
whose retained portions lie on a history graph for the unmodified RFDE.  The root is
not defined by a voltage threshold, a peak amplitude, or another
experimental observable: at the root the two retained complete histories
coincide.  For each fixed admissible \(\cP\), the root is unique near the
singular value and satisfies
\begin{equation}
 \label{eq:intro-root-law}
 \mu_{c,\cP}(\del,\eta)-\mu_{c,\cP}(\del,0)
 =\frac{K(\theta_0-\theta_1)}{4\alpha}\del^3\eta
  +O\bigl(\del^4\abs{\eta}+\del^3\eta^2\bigr),
 \qquad
 \alpha=\frac{\sqrt6}{4}.
\end{equation}
The remainder is uniform on the fixed parameter box and over bounded
admissible preparations with a common fixed logarithmic exponent \(p\).
Since \(4\alpha=\sqrt6\), the coefficient in
\eqref{eq:intro-root-law} is
\(K(\theta_0-\theta_1)/\sqrt6\).  In particular, when \(K>0\) and
\(\theta_0<\theta_1\), increasing \(\eta\) decreases the unscaled model
parameter \(\mu_c\) at leading order.  The exact finite-\(\del\) root may
depend on \(\cP\); the displayed displacement and its remainder are the
preparation-indexed statement proved here.

The conclusion is an identifiability result for critical-mode reductions.
At \(\eta=0\) the two delays enter the critical projection with weights
\(1/3\) and \(2/3\).  Varying \(\eta\) preserves these weights exactly and
also preserves the current-state total gain, yet it moves the root by order
\(\del^3\eta\) in the original unfolding parameter.  Consequently, even
retaining the full projected delay kernel is insufficient unless the
transverse delay response and its nonlinear return are controlled.  The
point is not merely that two networks with the same average delay can behave
differently; all projected delay tests agree in this family.  What differs
is how the two layers distribute the critical history into a stable mode.

Equation~\eqref{eq:intro-root-law} also separates the present contribution
from a new existence result for delayed canards.  The leading singular
canard is the classical planar one.  What is new in the analysis is the
remainder-controlled passage from a transverse, projection-invisible RFDE
perturbation to a simple complete-history connection root.  The proof must
retain the anisotropic transverse variables long enough to compute their
return, while controlling the selected histories on a domain that recedes
in the inner chart.  A fixed compact reduction supplies the local Taylor
coefficient but not, by itself, the root law: endpoint contributions at the
matching scale still have to be shown smaller than the claimed remainder.

Three analytical steps are responsible for \eqref{eq:intro-root-law}.  The
first is an anisotropic fold reduction in complete-history space.  The
critical voltage is of order \(\del\), whereas the transverse voltage and
transverse recovery variables occur at orders \(\del^2\) and \(\del^4\),
respectively.  In these coordinates the long physical delays become a fixed
history interval.  A Lyapunov--Perron construction gives a prepared invariant
history graph whose retained segments solve the uncut RFDE, together with an injective map from
its reduced planar dynamics to complete histories.  This avoids assigning a
backward evolution to an RFDE, whose natural solution operator is only a
semiflow; compare the phase-space foundations in
\citet{hale1993introduction}.

The second step extends the required mixed parameter estimates from a fixed
fold tube to the logarithmically growing interval
\[
 S_\del=\sqrt{2p\log(1/\del)}.
\]
This is long enough for Gaussian tail suppression to separate the local
coefficient from the prepared endpoints, but short enough that the graph
contraction and its \(C_\nu^1C_\eta^2\) jets remain uniform.  On this graph,
the transverse delay difference produces an order-\(\del^2\) return term in
the critical field.  Pairing that term with the adjoint Gaussian of the
singular canard yields the coefficient in \eqref{eq:intro-root-law}.

The third step turns this pairing into an actual connection theorem.  Along
the singular orbit, the planar variational equation has a tangent mode and a
normal mode with Gaussian growth at both ends.  We construct explicit
phase-normal Green operators on the attracting and repelling half-intervals:
the common phase removes the tangent mode, while opposite one-sided
projections remove the incoming normal mode.  A first integral of the
singular planar problem then defines a normalized gap.  Its
\(\nu\)-derivative is bounded away from zero, its \(\eta\)-derivative contains
the transverse-return coefficient, and the implicit-function theorem gives
\eqref{eq:intro-root-law}.  Finally, equality at the planar matching section,
uniqueness of the reduced flow, and injectivity of the history embedding
lift the zero gap to equality of the retained complete RFDE histories.

The theorem is local, and its preparation index is a genuine scope
qualification.  If a separately prescribed physical attracting and
repelling outer selection enters the local graph with parameter-coherent,
tame full-history boundary jets, the same Gaussian argument transfers
\eqref{eq:intro-root-law} to that selection.  We record this as a conditional
transfer statement in \cref{rem:physical-outer}, because those boundary estimates are not proved here for an
arbitrary RFDE Fenichel selection.  We likewise do not claim a global
spiking theorem.  General finite networks and the independent control of
frequency, amplitude, and a canard-safety coordinate require additional
transverse inverse and periodic-orbit results and are not part of the present
paper.

The paper is organized as follows.  In \Cref{sec:model-main} we introduce
the two-module RFDE, verify the exact projected-measure identities, and state
the main theorem.  \Cref{sec:history-graph} constructs the invariant
complete-history graph and establishes the mixed jets needed on the growing
fold region.  The reduced fold expansion and the transverse-return
coefficient are derived in \Cref{sec:fold-jet}.  In \Cref{sec:connection} we
construct the one-sided Green operators, define the normalized gap, and
prove the connection law.  \Cref{sec:numerics} gives an independent
method-of-steps sign and scale check, without using it as part of the proof.
The mathematical implications and the boundary with physical outer
selections and larger networks are discussed in \Cref{sec:discussion}.
